% ===== TITLE OPTIONS — uncomment one =====
% Option 1 (current favorite — the punchline):
%\newcommand{\papertitle}{Linear Methods Fail Where They Matter Most: \\
%Geometric Dissociation and Causal Subspaces in Neural Populations}
% Option 2 (provocative question):
%\newcommand{\papertitle}{Why Does Your Metric Disagree With Mine? \\
%Dimensionality Mediates Geometric Dissociation in the Brain}
% Option 3 (clean and direct):
%\newcommand{\papertitle}{The Metric Matters: Dimensionality Structure \\
%Determines Which Neural Similarity Measures Are Informative}
% Option 4 (causal framing forward):
%\newcommand{\papertitle}{Causal Subspaces Are Invisible to Linear Methods: \\
%Evidence from Brain-Wide Neuropixels Recordings}
% Option 5 (structured VAE angle):
%\newcommand{\papertitle}{Structured Disentanglement Reveals Causal Choice \\
%Pathways Hidden from Linear Discriminant Analysis}
% Option 6 (broadest — geometry + causality):
\newcommand{\papertitle}{Neural Geometry Is Not Metric-Neutral: \\
Dimensionality, Dissociation, and Causal Subspaces in Brain-Wide Neuropixels Recordings}
% Option 7 (shortest):
%\newcommand{\papertitle}{Geometric Dissociation in Neural Population Codes}
% Option 8 (silencing angle):
%\newcommand{\papertitle}{Optogenetic Ground Truth Validates Nonlinear \\
%Causal Subspaces Across the Mouse Brain}
% Option 9 (the finding, stated plainly):
%\newcommand{\papertitle}{A 3.4$\times$ Causal Signal That Linear Methods Miss: \\
%Nonlinear Subspace Discovery in Neural Populations}
% Option 10 (dual-finding):
%\newcommand{\papertitle}{Two Geometric Dissociations: Metrics Disagree \\
%About Neural Similarity, and Linear Methods Underestimate Causal Structure}

\documentclass{article}

% TMLR format
\usepackage[accepted]{tmlr}

\usepackage{amsmath,amssymb,amsfonts}
\usepackage{booktabs}
\usepackage{graphicx}
\usepackage{hyperref}
\usepackage{multirow}
\usepackage{xcolor}

% Placeholder command for pending results
\newcommand{\pending}[1]{\textcolor{red!70!black}{\textbf{[PENDING: #1]}}}

\title{\papertitle}

\author{
Elliot Tower \\
\texttt{elliot@elliottower.com}
}

\begin{document}
\maketitle

\begin{abstract}
We study how the choice of geometric similarity metric interacts with brain region dimensionality, and how this interaction conceals causal structure from standard linear methods. Using Neuropixels recordings across 73 brain regions in 10 mice \citep{steinmetz2019distributed}, we establish two geometric dissociations. \textbf{First}, linear (CKA) and nonlinear (UMAP Procrustes) similarity are strongly anti-correlated ($\rho = -0.85$, 95\% CI $[-0.87, -0.83]$), and this dissociation is fully mediated by effective dimensionality. \textbf{Second}, a structured variational autoencoder discovers causal choice subspaces that are 3.4$\times$ stronger than those found by linear discriminant analysis (LDA), winning in 73/73 brain regions ($p = 5.7 \times 10^{-14}$). This advantage generalizes across four task variables (choice, evidence strength, feedback, stimulus side; all $p < 10^{-13}$) and is validated against optogenetic silencing ground truth: LDA-identified subspaces are \emph{anti-correlated} with causal importance ($\rho = -0.73$, $p = 0.01$), while VAE subspaces show a positive trend, with the method difference significant (bootstrap $\Delta\rho = +1.06$, 95\% CI $[0.31, 1.51]$). The engagement and choice subspaces are near-orthogonal ($d_G = 1.82$), confirming that the VAE disentangles distinct cognitive variables. However, we also show that the VAE's interchange intervention accuracy (IIA) metric is vacuous for nonlinear methods---achieving comparable scores on random noise as on real data across both binary (IIA) and continuous (KL, JS, probability shift) metrics---motivating external validation over metric-internal evidence. A potent/null space decomposition confirms the mechanism: the geometric dissociation index is strongly anti-correlated with null-space fraction ($\rho = -0.86$, $p = 5 \times 10^{-22}$), meaning regions where metrics disagree most are precisely those with the largest choice-potent subspaces. The CKA--Procrustes dissociation replicates on the independent IBL Brain-Wide Map dataset ($\rho = -0.94$, 139 mice, 12 labs). These findings establish that linear subspace methods systematically fail in the highest-dimensional brain regions where causal signals are strongest, and that external causal validation (not IIA alone) is necessary to assess subspace quality.
\end{abstract}

\section{Introduction}

A growing body of work characterizes neural computation through the geometry of population activity \citep{cunningham2014dimensionality, gallego2017neural}. When $N$ neurons fire together, their joint activity traces trajectories in an $N$-dimensional space, and the low-dimensional structure of those trajectories reveals the underlying computation. Decision variables live in low-dimensional subspaces \citep{steinmetz2019distributed}. Motor plans occupy orthogonal manifolds \citep{gallego2017neural}. Representations drift over time within fixed-dimensional structures \citep{driscoll2022representational}.

This geometric view raises a methodological question with substantive consequences: \emph{does the choice of analysis method change the scientific conclusion?} The field's standard tools for identifying task-relevant subspaces---linear discriminant analysis (LDA), principal component analysis (PCA), and their variants---assume that the causal structure of neural population codes lies in linear subspaces. If this assumption fails, these methods will not merely be imprecise; they will systematically misidentify which brain regions carry the strongest causal signals.

We test this directly. Using 1,316 cross-session pairs across 73 brain regions in the Steinmetz et al.\ (2019) Neuropixels dataset, we find two geometric dissociations that together demonstrate the failure of linearity assumptions:

\textbf{Dissociation 1: Metrics disagree.} Linear (CKA) and nonlinear (UMAP Procrustes) similarity are strongly anti-correlated ($\rho = -0.85$). This is mediated by effective dimensionality: high-dimensional regions appear dissimilar to linear methods but highly similar to nonlinear ones. Once dimensionality is controlled for, the metrics agree.

\textbf{Dissociation 2: Linear methods miss the causal signal.} A structured variational autoencoder \citep{narayanaswamy2017learning} discovers choice subspaces that are 3.4$\times$ more causally effective than LDA subspaces, winning in all 73 brain regions. Critically, we validate this against optogenetic silencing data \citep{zatkahaas2021spatiotemporal}: LDA subspace quality is \emph{anti-correlated} with ground-truth causal importance ($\rho = -0.73$), meaning linear methods systematically identify the wrong regions as causally important. The VAE reverses this pattern.

\textbf{An honest caveat.} The interchange intervention accuracy (IIA) metric used to evaluate subspace quality is vacuous for nonlinear methods: both an unconstrained MLP and our structured VAE achieve high IIA on random Gaussian noise (mean $\approx 0.70$). This vacuity extends to continuous distributional metrics (KL divergence, JS divergence, probability shift), ruling out binary thresholding as the cause. Only linear DAS properly rejects noise (IIA $= 0.16$). This motivates our emphasis on external validation---optogenetic silencing ground truth, multi-task generalization, and engagement orthogonality---over metric-internal comparisons.

These findings converge on a single message: \emph{linear methods fail where they matter most}. The brain regions with the highest effective dimensionality, which are precisely the regions where causal choice signals are strongest, are the regions where LDA, CKA, and linear readouts are least informative. We interpret these results through the neural manifold framework \citep{gallego2017neural}: the choice subspace is the \emph{potent space} for decision output \citep{kaufman2014cortical}, and high-dimensional regions harbor large null spaces that confound linear methods while preserving manifold-level causal structure.

\section{Related Work}

\textbf{Neural manifolds and population-level computation.} \citet{gallego2017neural} argued that the fundamental unit of motor cortical computation is not the single neuron but the \emph{neural mode}---a population-level activation pattern spanning a low-dimensional manifold. Individual corticomotoneuronal cells need not encode any specific movement covariate; what matters is the trajectory through the space of modes. \citet{kaufman2014cortical} formalized this as the potent/null space distinction: preparatory activity lives in a null space orthogonal to the muscle-driving output dimensions, while execution activity enters the potent space. \citet{sadtler2014neural} demonstrated that subjects easily adapt to perturbations within the existing neural manifold but struggle with perturbations that require new modes, confirming that the manifold structure constrains behavior. Our evidence-choice subspace analysis directly tests this framework: the choice subspace is the potent space for decision output, and high-dimensional regions may harbor large null spaces that contain variance but no causal signal.

\textbf{Representational similarity.} CKA \citep{kornblith2019similarity} and RSA \citep{kriegeskorte2008representational} are the standard tools for comparing neural representations. CKA operates on kernel matrices and is therefore linear in the population activity; RSA abstracts from neuron identity by comparing representational dissimilarity matrices (RDMs), preserving relational structure across populations of different size \citep{kriegeskorte2008representational}. \citet{williams2021generalized} showed that many similarity measures are special cases of a generalized shape metric, but all within the linear/kernel family. Our Grassmannian analysis is complementary: RSA compares the global representational structure across all stimuli, while Grassmannian distance compares the subspace structure of the choice-relevant component specifically.

\textbf{Nonlinear methods.} UMAP \citep{mcinnes2018umap} and t-SNE \citep{vandermaaten2008visualizing} are widely used for visualization but less commonly for quantitative comparison. Procrustes analysis on UMAP embeddings has been used for cross-species comparison \citep{deitch2021representational}.

\textbf{Topological and dynamical methods.} Persistent homology has been applied to detect ring-like structures in head direction cells \citep{gardner2022toroidal} and grid cells \citep{rybakken2019decoding}. jPCA \citep{churchland2012neural} extracts rotational dynamics; DSA \citep{ostrow2023beyond} compares dynamical systems.

\textbf{Dimensionality and the cortical hierarchy.} Effective dimensionality scales unboundedly with neuron number \citep{stringer2024unbounded} and changes systematically along the cortical hierarchy \citep{huntenburg2024categorical}. \citet{harvey2024representational} proved that CKA measures average alignment between optimal linear readouts, with sensitivity depending on participation ratio---predicting our empirical finding that high-dimensional regions show low CKA even when manifold shape is conserved.

\textbf{Multimodal parcellation and convergent validity.} \citet{glasser2016multi} defined 180 cortical areas per hemisphere using converging evidence from architecture, function, connectivity, and topography. This multimodal convergence is the neuroscience gold standard for validating that region boundaries are biologically real. Our geometric clustering partially recovers anatomical organization (ARI $= 0.06$, NMI $= 0.30$), suggesting that our metrics capture functional rather than architectonic boundaries.

\textbf{Task transfer structure.} \citet{zamir2018taskonomy} built a computational taxonomy of visual tasks, showing that perceptual and geometric tasks cluster by computational type and that transfer structure is asymmetric. Our cross-region Grassmannian distance matrix is a neural analogue: regions whose choice subspaces are close on the Grassmannian have learned geometrically compatible representations, and the structure of this distance matrix may reveal a functional hierarchy from sensory to association to motor regions.

\textbf{Identifiable VAEs for neural data.} \citet{zhou2020pivae} proposed pi-VAE, which conditions the latent prior $p(\mathbf{z}|\mathbf{u})$ on task variables $\mathbf{u}$ via an exponential family parameterization, providing identifiability guarantees under mild conditions. Unlike our structured VAE, which separates $\mathbf{z}_{\text{choice}}$ from $\mathbf{z}_{\text{other}}$ via architecture and a classification head, pi-VAE achieves identifiability through the prior while using a single unsplit latent space. The semi-supervised disentanglement framework \citep{narayanaswamy2017learning} separates labeled task variables from unlabeled variance. The mechanistic independence condition \citep{khemakhem2020variational} provides identifiability guarantees when the generator Jacobian is sparser in factor-aligned coordinates. A natural hybrid combines pi-VAE's label-conditioned prior on $\mathbf{z}_{\text{choice}}$ with the structured split and classification head, gaining both identifiability and explicit subspace separation.

\textbf{Mechanistic interpretability of VAEs.} \citet{martinez2025vae_circuits} introduced causal effect strength (CES), intervention specificity, and circuit modularity as metrics for understanding how VAE encoders process inputs. Their key finding: FactorVAE develops highly specialized encoder channels with mediation values $> 0.7$ mapping directly to individual latent factors, while $\beta$-VAE distributes information across channels. Activation patching and causal mediation analysis---tools borrowed from transformer interpretability \citep{geiger2024causal}---can identify which encoder channels function as ``choice detectors'' in a structured VAE trained on neural data.

\textbf{IIA vacuity for nonlinear methods.} \citet{sutter2025iia} proved that unconstrained nonlinear alignment makes interchange intervention analysis (IIA) vacuous---any sufficiently flexible nonlinear map can achieve perfect IIA on arbitrary data. The linearity constraint of DAS \citep{geiger2024causal} is what makes IIA non-vacuous. Our structured VAE, despite having a structured latent split, retains enough encoder expressiveness to fit random noise, confirming Sutter et al.'s warning and motivating external validation.

\textbf{Dimension-level causal perturbation.} \citet{pereira2025residual} perturbed coding versus residual subspace dimensions in ALM, finding that residual dimensions causally drive choice through transient amplification. IIA \citep{geiger2024finding} and DAS \citep{geiger2024causal} perform formal causal interventions on learned subspaces in language models. Cross-region activation patching---swapping one region's choice-subspace output into another region's activity---is the biological analogue of path patching in transformer circuits.

\textbf{What is missing.} No prior work systematically compares metric families on the same neural data, validates subspace methods against optogenetic ground truth, tests whether the IIA vacuity concern applies to biological subspace estimation, or decomposes the potent/null space variance structure across dozens of brain regions simultaneously.

\section{Methods}

\subsection{Dataset}

We use the Steinmetz et al.\ (2019) dataset: Neuropixels recordings from 10 mice performing a two-alternative visual decision task. The dataset covers 73 brain regions with $\geq 15$ neurons per region (50 with $\geq 2$ recording sessions). We extract trial-averaged activity in a post-stimulus window (150--350ms) and compute binary choice labels (left vs.\ right). All cross-session pairs for the same region constitute the comparison set: 1,316 pairs across 50 multi-session regions.

\subsection{Geometric Invariants}

\textbf{Linear similarity (CKA).} Centered kernel alignment \citep{kornblith2019similarity} on trial $\times$ neuron activity matrices, using the linear kernel. We report both the standard (biased) estimator and the debiased HSIC estimator of \citet{song2012feature}, which corrects the finite-sample inflation identified by \citet{murphy2024biased} in the low-trial, high-neuron regime.

\textbf{Nonlinear manifold similarity (UMAP Procrustes).} 5-dimensional UMAP embedding \citep{mcinnes2018umap} with Procrustes alignment between sessions.

\textbf{Subspace distance (Grassmannian).} $k$-dimensional LDA subspaces with geodesic distance on $\operatorname{Gr}(k, n)$.

\textbf{Topological invariants.} Vietoris-Rips persistence diagrams via Ripser \citep{bauer2021ripser}.

\textbf{Dynamical invariants.} Simplified jPCA on choice-difference trajectories.

\textbf{Power-law exponent ($\alpha$).} Fit to PCA eigenvalue spectrum in log-log space (ranks 10--50), characterizing dimensionality structure.

\subsection{Structured VAE for Subspace Discovery}

We train a structured variational autoencoder \citep{narayanaswamy2017learning} per region. The encoder maps population activity to two latent groups: $\mathbf{z}_{\text{choice}}$ ($k$ dimensions, $k \in \{1,2,3\}$) trained with an auxiliary classification loss to predict left/right choice, and $\mathbf{z}_{\text{other}}$ ($m = 15$ dimensions) capturing remaining variance. Both groups jointly reconstruct the population activity via the ELBO. Training: 300 epochs, $\alpha_{\text{task}} = 10$, $\beta_{\text{KL}} = 1$, across all 39 sessions.

The encoder weight matrix for $\mathbf{z}_{\text{choice}}$ directly provides the choice subspace directions, replacing the LDA+PCA pipeline.

\subsection{Interchange Intervention Analysis (IIA)}

For each region, we swap the choice-subspace projection between opposite-evidence trial pairs and measure how often a choice classifier flips its prediction. IIA $= 1$ means the subspace is perfectly causal; IIA $= 0.5$ means random. We compute IIA for both the VAE and LDA subspaces, and against a random-direction null.

\subsection{Optogenetic Silencing Validation}

We validate subspace quality against ground-truth causal importance from the Zatka-Haas et al.\ (2021) optogenetic inactivation dataset: 52 cortical coordinates, 47,002 laser trials across 10 mice. For each silencing coordinate mapped to a Steinmetz brain region ($n = 12$ direct matches, $n = 16$ with hierarchical Allen CCF grouping), we compute the behavioral effect size (Euclidean distance in choice-probability space between laser-on and laser-off trials on matched-contrast conditions). The Spearman correlation between silencing effect and subspace IIA provides external validation independent of any metric-internal evaluation.

\subsection{Multi-Task Generalization}

Beyond the primary choice variable, we test subspace discovery on three additional task variables: evidence strength (high vs.\ low contrast), feedback (correct vs.\ error), and stimulus side (left vs.\ right). For each variable, we train a separate structured VAE and compute IIA against the corresponding LDA baseline.

\subsection{Engagement Orthogonality}

We derive an engagement variable from behavioral signatures: disengaged trials (NoGo responses on easy stimuli where the correct response is obvious) vs.\ engaged trials (correct responses on matched stimuli). We estimate engagement and choice subspaces in a pre-stimulus window (0--150ms) and post-stimulus window (250--450ms) respectively, then compute their Grassmannian distance. Near-orthogonality ($d_G \approx \pi/2$ per dimension) would confirm that the VAE disentangles genuinely distinct cognitive variables rather than capturing a single dominant signal.

\section{Results}

\subsection{Dissociation 1: Linear and Nonlinear Metrics Anti-Correlate}

Across 1,316 cross-session pairs and 50 brain regions, CKA and UMAP Procrustes are strongly anti-correlated: Spearman $\rho = -0.85$, 95\% CI $[-0.87, -0.83]$. Motor cortex (MOs) exemplifies the extreme: CKA $= 0.07$ (near-zero linear similarity) but UMAP Procrustes $= 0.98$ (near-maximal nonlinear similarity).

\textbf{The dissociation is robust.} Bootstrap 95\% CI: $[-0.867, -0.830]$. Leave-one-mouse-out jackknife: $\rho \in [-0.872, -0.835]$. Leave-one-region-out: $\rho \in [-0.863, -0.841]$. Neuron sub-sampling: $\rho = -0.81$ ($p < 10^{-7}$). Spectral denoising (Marchenko--Pastur thresholding) \emph{strengthens} the result: $\rho = -0.889$.

\textbf{Debiased CKA confirms the dissociation.} \citet{murphy2024biased} showed that standard CKA inflates similarity in the low-trial, high-neuron regime. Using the unbiased HSIC estimator \citep{song2012feature}, the CKA--Procrustes correlation drops from $\rho = 0.85$ (biased) to $\rho = 0.45$ (debiased), with partial correlation controlling for neuron count of $0.34$ (Appendix~\ref{app:debiased_cka}). The metric dissociation is \emph{real}---not an artifact of biased CKA estimation---though the biased estimator overstates its strength.

\textbf{Dimensionality is the mechanism.} Partial correlation controlling for power-law exponent $\alpha$ reverses the sign: partial $r = +0.44$. Controlling for neuron count alone gives partial $r = +0.13$. Once dimensionality is accounted for, CKA and Procrustes \emph{agree}. This empirically validates the theoretical prediction of \citet{harvey2024representational}.

\textbf{The anti-correlation is metric-specific.} Sliced Wasserstein distance is uncorrelated with dimensionality ($\rho = 0.02$, $p = 0.47$), while SPD Riemannian distance tracks it ($\rho = 0.44$). The dissociation is not a generic ``all metrics disagree'' phenomenon but a specific interaction between kernel-based and manifold-based methods mediated by spectral structure.

\subsection{Spectral Universality and Dynamic Dimensionality}

Singular value spectra are highly conserved across animals (mean $r = 0.94 \pm 0.06$), but temporal modes are idiosyncratic (mean cosine $= 0.09$). The dimensionality \emph{structure} is universal; the \emph{content} is animal-specific.

Dimensionality is not fixed. Under hard task conditions (low contrast), $\alpha$ increases by $+12.4$ (Wilcoxon $p = 3.0 \times 10^{-7}$) while manifold shape is preserved (Procrustes $= 0.90$). Dimensionality compression is an active computational strategy, not a fixed anatomical property.

Rotational dynamics appear in 72/73 regions above a temporal-shuffle null, but strength varies 50$\times$ (range $0.008$--$0.40$). Persistent homology reveals almost no topological structure: only 9/73 regions have $\beta_1 > 0$.

\subsection{Evidence-Choice Subspace Alignment}

For each region, we estimate evidence and choice subspaces and measure their Grassmannian distance. Two strategies emerge:

\begin{enumerate}
\item \textbf{Low-dimensional regions} maintain stable evidence-choice alignment across conditions (high overlap, small choice subspace shift under evidence change).
\item \textbf{High-dimensional regions} show large choice subspace rotations when evidence is experimentally changed ($\rho = -0.59$, permutation $p < 10^{-6}$), maintaining separate, flexibly rotatable representations.
\end{enumerate}

Cross-condition (easy vs.\ hard) subspace rotation also anti-correlates with $\alpha$ ($\rho = -0.65$, 95\% CI $[-0.76, -0.49]$, permutation $p < 0.002$).

\subsection{Dissociation 2: Learned Subspaces Amplify Causal Signal}

\textbf{VAE subspaces are 3.4$\times$ more causal than LDA.} The VAE produces mean IIA $= 0.56$ (SD $0.08$) vs.\ LDA mean IIA $= 0.17$ (SD $0.09$), with Cohen's $d = 4.5$. The VAE wins in 73/73 regions (Wilcoxon $W = 2701$, $p = 5.7 \times 10^{-14}$). It exceeds the random-direction null in all 73 regions (mean null $= 0.23$).

\textbf{Choice subspace dimensionality ablation.} The causal signal is effectively one-dimensional: even $k = 1$ achieves a 3.7$\times$ improvement over LDA (72/73 wins), and performance is stable across $k \in \{1, 2, 3\}$.

\begin{table}[h]
\centering
\small
\caption{Effect of choice subspace dimensionality $k$ on structured VAE IIA across 73 brain regions.}
\label{tab:vae_k_ablation}
\begin{tabular}{lccccc}
\toprule
$k$ & VAE IIA & LDA IIA & Ratio & VAE wins & Choice acc. \\
\midrule
1 & $0.548 \pm 0.093$ & $0.150 \pm 0.091$ & 3.7$\times$ & 72/73 & 0.99 \\
2 & $0.555 \pm 0.088$ & $0.159 \pm 0.097$ & 3.5$\times$ & 73/73 & 0.98 \\
3 & $0.561 \pm 0.080$ & $0.166 \pm 0.094$ & 3.4$\times$ & 73/73 & 0.99 \\
\bottomrule
\end{tabular}
\end{table}

\textbf{Low-dimensional regions benefit most.} Low-$\alpha$ (high-dimensional) regions show a larger VAE improvement ($\Delta$IIA $= 0.46$) than high-$\alpha$ regions ($\Delta$IIA $= 0.33$). VAE IIA anti-correlates with $\alpha$ ($\rho = -0.48$, $p = 1.5 \times 10^{-5}$), reversing the pattern seen with LDA.

\textbf{VAE finds genuinely different subspaces.} The VAE and LDA subspaces are nearly orthogonal: mean Grassmannian distance $= 2.34$ (SD $0.15$), far exceeding $\pi/2 \approx 1.57$ per dimension. The VAE does not refine the LDA solution---it finds a fundamentally different set of directions.

\textbf{Model confidence tracks causal effectiveness.} Posterior uncertainty anti-correlates with IIA ($\rho = -0.26$, $p = 0.028$): regions where the model is most confident are those where interventions are most effective.

\subsection{The IIA Vacuity Problem}
\label{sec:vacuity}

Before interpreting the 3.4$\times$ improvement as evidence that the VAE finds better causal subspaces, we must address a fundamental concern raised by \citet{sutter2025iia}: is IIA a valid metric for nonlinear methods?

We test three methods on real neural data versus random Gaussian noise with matched labels:

\begin{table}[h]
\centering
\small
\caption{IIA on real neural data vs.\ random Gaussian noise (73 regions). A non-vacuous method should show high IIA on real data and chance IIA ($\approx 0.5$) on noise.}
\label{tab:sutter}
\begin{tabular}{lccl}
\toprule
Method & Real IIA & Random IIA & Interpretation \\
\midrule
Unconstrained MLP & --- & $0.70 \pm 0.11$ & Vacuous (fits noise) \\
Structured VAE & $0.69 \pm 0.11$ & $0.70 \pm 0.11$ & Also vacuous on noise \\
Linear DAS & --- & $0.16 \pm 0.09$ & Non-vacuous (rejects noise) \\
\bottomrule
\end{tabular}
\end{table}

The structured VAE achieves comparable IIA on random noise ($0.70$) as on real data ($0.69$). The structured latent split alone does not prevent the encoder from fitting arbitrary labels. This means \emph{the 3.4$\times$ IIA improvement over LDA could in principle reflect encoder expressiveness rather than genuine causal subspace discovery.}

One might hypothesize that this vacuity is an artifact of the binary IIA metric itself---perhaps continuous measures of distributional shift would distinguish real from random data even when the binary flip rate does not. We test this directly by computing four continuous metrics alongside IIA: KL divergence ($D_{KL}(p_{\text{orig}} \| p_{\text{swap}})$), Jensen-Shannon divergence, probability shift ($|\Delta P(\text{choice})|$), and logit difference ($|\Delta \ell|$). Results across all 73 regions:

\begin{table}[h]
\centering
\small
\caption{Continuous metrics on real vs.\ random data (Exp72). If vacuity were a property of the binary metric, continuous metrics would show real $\gg$ random. They do not.}
\label{tab:continuous}
\begin{tabular}{llccl}
\toprule
Method & Metric & Real & Random & Wilcoxon $p$ \\
\midrule
\multirow{4}{*}{VAE} & IIA & 0.688 & 0.695 & 0.97 \\
 & KL divergence & 2.51 & 2.59 & 0.99 \\
 & JS divergence & 0.397 & 0.408 & 0.99 \\
 & Prob.\ shift & 0.655 & 0.666 & 0.98 \\
\midrule
\multirow{2}{*}{MLP} & IIA & 0.687 & 0.695 & 0.95 \\
 & Logit diff.\ & 8.93 & 8.13 & $3.3 \times 10^{-10}$ \\
\midrule
Linear DAS & IIA & 0.167 & 0.163 & 0.17 \\
\bottomrule
\end{tabular}
\end{table}

The VAE shows real $\approx$ random on \emph{every} continuous metric ($p > 0.95$ in all cases). The vacuity is not an artifact of binary thresholding---it reflects genuine encoder flexibility. The one exception, MLP logit difference ($p = 3.3 \times 10^{-10}$), reflects the unconstrained encoder producing larger raw logit magnitudes on real data despite identical binary flip rates, suggesting scale differences without discrimination improvement.

This result motivates two responses: (1) external validation that does not depend on IIA or any intervention metric, and (2) ablation experiments that isolate the source of the improvement.

\textbf{Shuffled-label control.} To test whether the VAE learns genuine choice structure despite being able to fit noise, we trained with randomly shuffled choice labels on real neural data (5 shuffles per region). Real IIA (mean $= 0.625$) significantly exceeds shuffled-label IIA (mean $= 0.430$) in all 24 tested regions (Wilcoxon $W = 300$, $p = 6.0 \times 10^{-8}$). Shuffled-evidence IIA is even lower (mean $= 0.424$), confirming that the VAE exploits genuine statistical structure in the choice labels. The encoder \emph{can} fit noise, but it finds stronger structure in real data.

\textbf{Linear VAE ablation.} To isolate whether the benefit comes from encoder nonlinearity or from the structured latent split, we compared nonlinear and linear structured VAEs. The nonlinearity effect is negligible: mean IIA difference $= 0.001$, Wilcoxon $p = 0.65$. Linear and nonlinear encoders achieve equivalent IIA when both have the structured split and classification loss. This confirms that the structured architecture, not nonlinear expressiveness, drives the improvement---and explains why the VAE matches noise performance: any sufficiently flexible encoder (linear or nonlinear) can fit random labels, but the structured loss channels real data into the choice subspace.

\subsection{External Validation: Optogenetic Silencing Ground Truth}

The strongest evidence for the VAE's superiority comes not from IIA comparisons but from external validation against optogenetic silencing data---a test that is completely independent of the IIA metric.

Across 12 brain regions matched between Steinmetz and Zatka-Haas et al.\ (2021):

\begin{table}[h]
\centering
\small
\caption{Spearman correlation between subspace quality and optogenetic silencing effect (ground-truth causal importance). Bootstrap BCa CIs from 10,000 resamples.}
\label{tab:silencing}
\begin{tabular}{lccc}
\toprule
Metric & $\rho$ vs.\ silencing & $p$ & 95\% CI \\
\midrule
LDA IIA & $-0.73$ & $0.01$ & $[-0.93, -0.27]$ \\
VAE IIA & $+0.33$ & $0.30$ & --- \\
$\Delta\rho$ (VAE $-$ LDA) & $+1.06$ & $< 0.005$ & $[0.31, 1.51]$ \\
\bottomrule
\end{tabular}
\end{table}

\textbf{LDA subspaces are anti-correlated with causal importance} ($\rho = -0.73$, permutation $p = 0.01$). Regions where optogenetic silencing most disrupts behavior---the ground-truth causally important regions---are the regions where LDA-based IIA is \emph{lowest}. Linear methods systematically identify the wrong regions as causally important.

\textbf{The VAE reverses this pattern.} The VAE correlation is positive ($\rho = +0.33$), though not individually significant at $n = 12$. The critical test is the \emph{difference}: $\Delta\rho = +1.06$, 95\% bootstrap CI $[0.31, 1.51]$, permutation $p < 0.005$. The CI excludes zero, confirming that the VAE provides a significantly better match to optogenetic ground truth than LDA.

\textbf{Hierarchical region expansion.} Using Allen CCF ontology to group Steinmetz sub-regions under parent silencing areas expands the matched set from 12 to 16 regions. The expanded analysis shows VAE $\rho = +0.42$ (perm $p = 0.11$, trending), LDA $\rho = -0.29$, and the delta-rho CI still excludes zero ($\Delta\rho = 0.87$, 95\% CI $[0.09, 1.43]$).

Per-mouse optogenetic validation is reported in Appendix~\ref{app:per_mouse}.

\subsection{Multi-Task Generalization}

The VAE advantage is not specific to the choice variable. Across three additional task variables, the VAE wins 73/73 regions for every variable:

\begin{table}[h]
\centering
\small
\caption{VAE vs.\ LDA IIA across four task variables. All 73 regions, all 39 sessions.}
\label{tab:multitask}
\begin{tabular}{lcccc}
\toprule
Task variable & VAE IIA & LDA IIA & VAE wins & Wilcoxon $p$ \\
\midrule
Choice & $0.56 \pm 0.08$ & $0.17 \pm 0.09$ & 73/73 & $5.7 \times 10^{-14}$ \\
Evidence strength & $0.64 \pm 0.13$ & $0.26 \pm 0.10$ & 73/73 & $5.7 \times 10^{-14}$ \\
Feedback & $0.58 \pm 0.11$ & $0.21 \pm 0.09$ & 73/73 & $5.7 \times 10^{-14}$ \\
Stimulus side & $0.66 \pm 0.11$ & $0.29 \pm 0.10$ & 73/73 & $5.7 \times 10^{-14}$ \\
\bottomrule
\end{tabular}
\end{table}

The improvement is largest for evidence strength (2.4$\times$) and smallest for stimulus side (2.3$\times$), but all four show overwhelming VAE dominance. The anti-correlation between $\alpha$ and VAE IIA holds for all task variables: evidence strength ($\rho = -0.75$, $p < 10^{-13}$), feedback ($\rho = -0.70$, $p < 10^{-11}$), stimulus side ($\rho = -0.76$, $p < 10^{-14}$).

\subsection{Engagement-Choice Orthogonality}

The VAE identifies engagement and choice subspaces that are near-orthogonal across all 59 regions with sufficient data:

\begin{itemize}
\item \textbf{Grassmannian distance}: $d_G = 1.82$ (LDA), $d_G = 2.41$ (VAE). Both exceed $\pi/2 \approx 1.57$, confirming orthogonality.
\item \textbf{VAE wins 59/59 regions} for engagement-choice separation.
\item \textbf{Cross-variable IIA}: swapping the engagement subspace between opposite-choice trials produces mean IIA $= 0.17$---near-random, confirming that engagement interventions do not affect choice and vice versa.
\end{itemize}

This demonstrates that the VAE is not merely finding a single dominant nonlinear direction that correlates with everything. It discovers geometrically orthogonal subspaces for genuinely distinct cognitive variables.

\subsection{Cross-Task Subspace Geometry}

We compute pairwise Grassmannian distances between all four task-variable subspaces per region. Two patterns emerge:

\textbf{Choice and stimulus side share structure.} LDA Grassmannian distance between choice and stimulus side is $d_G = 0.585$ (the closest pair), consistent with stimulus side directly driving choice. All other task pairs are near-orthogonal ($d_G \approx 1.3$).

\textbf{VAE separates all task variables more strongly.} The VAE produces larger pairwise distances than LDA in 73/73 regions ($p < 10^{-14}$), with mean $d_G = 2.36$ (VAE) vs.\ $1.20$ (LDA). The VAE finds subspaces that are more orthogonal---better disentangled---than LDA's projections.

\textbf{Multiplexing varies by region.} The most multiplexed regions (lowest mean pairwise distance, most overlapping task subspaces) include OT, LH, and VISa---regions at the interface of multiple processing streams. The least multiplexed (most orthogonal) include PIR, EPd, and AUD---unimodal sensory regions.

\textbf{Temporal emergence of the choice subspace.} To test whether the choice subspace reflects a causal decision signal rather than a post-hoc motor artifact, we computed IIA in 50ms sliding windows from stimulus onset through response. Both LDA and VAE subspaces show pre-response emergence: 100\% of regions with detectable onset show rising IIA before the behavioral response. Mean onset occurs at bin 35.9 ($\approx 360$ms post-stimulus), consistent with the time course of decision formation in frontal and parietal cortex.

\subsection{Cross-Dataset Replication: IBL Brain-Wide Map}
\label{sec:ibl}

To test whether the CKA--Procrustes dissociation generalizes beyond a single dataset, we replicated the analysis on the IBL Brain-Wide Map \citep{ibl2025brainwide}, an independent multi-laboratory dataset comprising 139 mice across 12 institutions. After matching brain regions between datasets (requiring $\geq 3$ sessions and $\geq 15$ neurons per region), 11 regions were available for direct comparison: CA1, CA3, DG, GPe, LP, POST, RSPd, SUB, VISam, VISl, and VISp.

The CKA--Procrustes anti-correlation replicates strongly in the IBL data ($\rho = -0.94$, $p = 5.5 \times 10^{-5}$; permutation test $p < 10^{-4}$, $n = 10{,}000$ permutations). The effect size is comparable to the original Steinmetz result ($\rho = -0.90$ in the matched subset, $-0.85$ across all 50 regions), despite substantial differences in recording pipeline, spike sorting, and session protocols across the 12 IBL labs. Error bars ($\pm 1$ SE across sessions per region) are modest, confirming that within-region variability does not drive the correlation.

Geometric type classification transfers for 7 of 11 matched regions (64\%). The four mismatches (CA3, DG, RSPd, VISl) are all hippocampal or higher visual regions near the CKA--Procrustes boundary, where small shifts in absolute metric values can flip the binary classification while preserving the overall anti-correlation.

Absolute power-law exponents do \emph{not} transfer across datasets ($\rho = -0.27$, n.s.), indicating that the \emph{relative} geometric relationship between metrics is robust while absolute dimensionality estimates are sensitive to recording pipeline differences (neuron yield, spike sorting, session count). This dissociation between transferable relationships and non-transferable absolutes is itself informative: it suggests that the CKA--Procrustes anti-correlation reflects a genuine property of neural population geometry rather than an artifact of any particular recording setup.

\section{Discussion}

\subsection{Linear Methods Fail Where They Matter Most}

Our central claim is not merely that nonlinear methods outperform linear ones---this has been shown before in various settings. The novel finding is the \emph{interaction} with dimensionality: linear methods fail most severely in the highest-dimensional brain regions, and these are precisely the regions where optogenetic silencing reveals the strongest causal role in behavior.

This creates a systematic blind spot. A researcher using LDA to identify causally important regions would conclude that low-dimensional sensory regions carry the strongest choice signals (because LDA IIA is highest there), while high-dimensional association and hippocampal regions are causally weak. The optogenetic ground truth shows the opposite: the high-dimensional regions are the most causally important, but their causal subspaces are invisible to linear methods.

\subsection{The Neural Manifold Hypothesis as Unifying Framework}

Our results are naturally interpreted through the neural manifold framework of \citet{gallego2017neural}. The evidence subspace identified by the VAE corresponds to the \emph{potent space} for choice output \citep{kaufman2014cortical}---the subspace of population activity that causally drives the downstream binary decision. The remaining variance, which may be high-dimensional in association regions like MOs, lives in the \emph{null space}---it is present, it contributes to CKA dissimilarity, but it does not drive behavior.

This potent/null decomposition explains the CKA--Procrustes dissociation. A region like MOs shows CKA $\approx 0$ because the high-dimensional null space dominates the kernel matrix, making sessions appear dissimilar. But UMAP Procrustes $\approx 1$ because the manifold \emph{shape}---including the low-dimensional choice-potent subspace embedded within it---is conserved. CKA tests whether two regions share the same \emph{task-specific coordinates} (same directions of variance); Procrustes tests whether they share the same \emph{universal manifold structure} (same geometric shape). High Procrustes with low CKA means: same universal manifold, different subspace sampled within it.

The within-manifold / outside-manifold learning distinction of \citet{sadtler2014neural} makes a testable prediction: if the VAE's choice subspace represents a within-manifold representation, then IIA interventions should be robust to noise (the manifold provides a scaffold for downstream interpretation). Conversely, if choice encoding is outside-manifold for some regions, IIA should be noisier. This predicts a correlation between manifold dimensionality and IIA reliability across regions.

We tested this framework directly by decomposing each region's activity into its choice subspace (potent) and complement (null), computing the fraction of total variance in each. Across 73 regions, LDA places 99.8\% of variance in the null space on average, while the VAE places 84.3\%---consistent with the VAE discovering a larger potent subspace that captures more of the choice-relevant structure. Despite this difference, both methods decode choice from the potent subspace with comparable accuracy ($\sim$66\%), confirming that the potent subspace---however small---carries the causal signal.

Critically, the geometric dissociation index $\alpha$ is strongly anti-correlated with the null-space fraction: $\rho = -0.86$ ($p = 5 \times 10^{-22}$) for the VAE and $\rho = -0.68$ ($p = 3 \times 10^{-11}$) for LDA. Regions where linear and nonlinear metrics disagree most are precisely the regions with the smallest null-space fractions---i.e., the regions where the most variance is concentrated in the choice-potent subspace. Conversely, potent-space decoding accuracy does not correlate with $\alpha$ ($\rho = -0.07$, $p = 0.57$ for VAE; $\rho = -0.03$, $p = 0.77$ for LDA), confirming that the \emph{informativeness} of the potent subspace is constant across regions while its \emph{relative size} drives the metric dissociation. This is exactly the prediction of the potent/null framework: high-dimensional regions have large null spaces that inflate CKA dissimilarity while preserving manifold-level causal structure.

\subsection{The IIA Vacuity Problem and External Validation}

We take the Sutter dilemma seriously. Our structured VAE achieves IIA $= 0.70$ on random Gaussian noise---the same as an unconstrained MLP. This means IIA alone cannot distinguish genuine causal subspace discovery from encoder flexibility. The structured latent split (z\_choice/z\_other) does not provide sufficient constraint to prevent overfitting.

Three lines of evidence nevertheless support the VAE's validity:

\begin{enumerate}
\item \textbf{Optogenetic validation} ($\S$5.6): The VAE-silencing correlation is positive while LDA's is strongly negative, and the difference is significant. This test is completely independent of IIA.
\item \textbf{Cross-task disentanglement} ($\S$5.8): The VAE discovers near-orthogonal subspaces for engagement vs.\ choice, with cross-variable IIA at chance ($0.17$). A vacuous method would not produce orthogonality between genuinely distinct cognitive variables.
\item \textbf{Multi-task generalization} ($\S$5.7): The VAE advantage holds for all four task variables with the same dimensionality interaction pattern. A single overfitting mechanism would not generalize this consistently.
\end{enumerate}

The shuffled-label control ($\S$5.4) confirms this: real IIA significantly exceeds shuffled-label IIA in all 24 regions ($p = 6 \times 10^{-8}$), establishing that the VAE learns genuine choice structure despite also being able to fit noise. The linear VAE ablation further shows that the nonlinearity itself contributes nothing (IIA difference $= 0.001$, $p = 0.65$); the structured split and classification loss are the active ingredients.

\subsection{Dimensionality as Computational Strategy}

The task-difficulty modulation ($\Delta\alpha = +12.4$, $p = 3 \times 10^{-7}$) while preserving manifold shape demonstrates that dimensionality is an active computational variable, not a fixed property. Under difficult conditions, circuits compress into fewer dimensions---concentrating variance along task-relevant axes. The evidence-choice subspace analysis reveals two strategies: low-dimensional regions route through shared subspaces (efficient, inflexible), while high-dimensional regions maintain separate, rotatable representations (flexible, requiring nonlinear readout).

\subsection{Toward Neural Taskonomy and RSA of Subspaces}

Our pairwise Grassmannian distances between regions' choice subspaces define a \emph{Grassmannian dissimilarity matrix} (GDM)---an analogue of the representational dissimilarity matrix (RDM) used in RSA \citep{kriegeskorte2008representational}, but operating on subspaces rather than activity patterns. Where RSA asks ``do two conditions evoke similar population responses?'', the GDM asks ``do two regions use similar geometric strategies to encode choice?'' Correlating the GDM with CKA-based and Procrustes-based RDMs, anatomical distance, and functional connectivity would reveal whether subspace similarity captures biologically meaningful organization that standard metrics miss.

This extends naturally to a \emph{neural taskonomy} \citep{zamir2018taskonomy}. Just as Zamir et al.\ discovered that visual tasks form a computational taxonomy with asymmetric transfer, the GDM may reveal a functional hierarchy in which sensory regions' subspaces are close to each other but distant from motor regions', with association regions in between. The structure of this graph---and whether it predicts which region's lesion most disrupts another region's decoding---is a direct test of the manifold framework applied to brain-wide circuit organization.

We constructed such a GDM by computing per-session Grassmannian distances between all co-recorded region pairs (projecting into a shared 30-dimensional PCA ambient space per session), then averaging across sessions. Coverage spans 719 of 2,628 possible region pairs (mean 1.9 sessions per pair).

Two striking patterns emerge. \textbf{First}, the LDA and VAE Grassmannian distance matrices are nearly identical (Mantel $\rho = 0.97$, $p < 10^{-4}$). Despite finding very different choice subspaces---LDA captures a 1D discriminant while the VAE finds a 3D nonlinear manifold---the \emph{relative geometry} between regions is conserved. Both methods agree on which regions use similar geometric strategies to encode choice. \textbf{Second}, both GDMs are strongly anti-correlated with CKA similarity ($\rho = -0.96$ for both LDA and VAE). Regions that appear similar in CKA's kernel space have \emph{dissimilar} choice subspaces, reinforcing the dissociation: CKA captures shared variance structure (dominated by null-space dimensions) while the GDM captures shared choice geometry (the potent subspace). Neither GDM recovers traditional functional groupings (Appendix~\ref{app:gdm_hierarchy})---within-group Grassmannian distances are not smaller than between-group distances---suggesting that choice subspace geometry is organized along a dimension orthogonal to the classical sensory--association--motor hierarchy.

\subsection{Identifiable Subspace Discovery via pi-VAE}

The IIA vacuity problem (\S\ref{sec:vacuity}) suggests that our structured VAE's encoder may have too much expressiveness---it can fit arbitrary labels. A principled remedy is to constrain the model with identifiability guarantees. pi-VAE \citep{zhou2020pivae} achieves this by conditioning the latent prior on task variables: $p(\mathbf{z}|\mathbf{u}) = \prod_i \frac{Q_i(z_i)}{Z_i(\mathbf{u})} \exp\left[\sum_j T_j(z_i) \lambda_{i,j}(\mathbf{u})\right]$, where $\mathbf{u}$ is the choice label. Under mild conditions (injective decoder, differentiable sufficient statistics, enough distinct label values), the latent representation is identifiable up to a linear transformation.

We propose a hybrid ``pi-structured-VAE'' that combines both approaches: the label-conditioned prior is applied \emph{only} to $\mathbf{z}_{\text{choice}}$, while $\mathbf{z}_{\text{other}}$ retains a standard Gaussian prior. This preserves the structured split (explicit choice subspace) while gaining identifiability on the choice-relevant dimensions. The classification head provides an additional signal beyond the prior.

The linear VAE ablation ($\S$5.4) partially addresses this: the linear structured VAE achieves equivalent IIA to the nonlinear version (mean difference $= 0.001$, $p = 0.65$), suggesting that the structured split, not encoder expressiveness, is the active ingredient. A full pi-structured-VAE hybrid ablation (Appendix~\ref{app:pivae}) shows the label-conditioned prior does not improve and actually disrupts the structured split (hybrid IIA $= 0.036$ vs.\ structured VAE $= 0.941$).

\subsection{Cross-Region Activation Patching}

The biological analogue of path patching in transformer circuits \citep{geiger2024finding} is cross-region activation patching: replacing one region's choice-subspace output with another region's and measuring whether downstream decoding changes. For simultaneously recorded region pairs $(A, B)$, we project $A$'s activity onto its VAE choice subspace, transplant this component into $B$'s activity (replacing $B$'s own projection onto $A$'s subspace), and measure whether a classifier trained on $B$'s activity flips its choice prediction.

If patching $A \to B$ flips $B$'s decoded choice, then $A$'s choice subspace lies in $B$'s potent space---the two regions share causal geometry. The directed patching graph across all 73 regions would reveal information flow: regions with high \emph{outgoing} patching strength are causal sources (their choice signal propagates), while regions with high \emph{incoming} strength are receivers (their choice decoding depends on others). The asymmetry of this graph tests the prediction that information flows from sensory to decision to motor regions.

Cross-region activation patching across 1,438 directed pairs reveals directional information flow: thalamic regions (MG, MD) are net causal sources, while visual and striatal regions are net receivers (Appendix~\ref{app:patching}).

\subsection{Mechanistic Interpretability of the VAE Encoder}

Beyond using the VAE as a tool for subspace discovery, we can ask how the encoder \emph{itself} processes neural activity---applying mechanistic interpretability to the analysis method rather than the brain. Following the multi-level causal intervention framework of \citet{martinez2025vae_circuits}, we compute three metrics per trained VAE: causal effect strength (CES $= \mathbb{E}[\|\mathbf{D}(\mathbf{z}) - \mathbf{D}(\tilde{\mathbf{z}}_i)\|_2]$, measuring how much intervening on latent dimension $i$ changes the reconstruction), intervention specificity ($S(i) = 1/(H(\mathbf{p}_i) + \epsilon)$, measuring whether changes are localized), and circuit modularity ($M = 1 - \frac{1}{\binom{K}{2}}\sum_{i<j}|\rho(\Delta \mathbf{a}_i, \Delta \mathbf{a}_j)|$, measuring whether different latent dimensions use separate encoder pathways).

If the structured VAE develops ``choice detector'' channels in its encoder---analogous to FactorVAE's specialized shape-detecting channels \citep{martinez2025vae_circuits}---this would confirm that the choice/other split is not merely a labeling convenience but reflects a genuine architectural specialization learned from the data.

Mechanistic analysis of the VAE encoder reveals choice dimensions carry $2.73\times$ more causal effect strength than other dimensions, with $2.72\times$ lift over random splits (Appendix~\ref{app:circuits}).

\subsection{Practical Implications}

For neural data analysis, our findings imply a simple diagnostic: compute the power-law exponent $\alpha$ before selecting analysis methods. For $\alpha > 3$ (low-dimensional), linear methods (CKA, LDA) are informative. For $\alpha < 1$ (high-dimensional), nonlinear methods are necessary---and linear methods may be actively misleading. For intermediate regimes, both should be reported.

For the interpretability community, the IIA vacuity result is a cautionary finding. IIA is widely used to evaluate causal subspace quality in language models; our results suggest that external validation (behavioral perturbation, out-of-distribution generalization) is necessary alongside metric-internal evaluation, particularly for nonlinear methods.

\subsection{Limitations}

\textbf{Dataset scale.} The Steinmetz dataset comprises 10 mice and 73 regions. The IBL replication (\S\ref{sec:ibl}) confirms the core anti-correlation on an independent dataset (139 mice, 12 labs, 11 matched regions), but absolute power-law exponents do not transfer ($\rho = -0.27$, n.s.), limiting cross-dataset dimensionality comparisons. Per-region significance testing would require denser coverage than the current 3--5 sessions per region in IBL.

\textbf{Silencing sample size.} $n = 12$ direct matches ($n = 16$ with hierarchical grouping) for the optogenetic validation. Power analysis shows the LDA anti-correlation ($\rho = -0.73$) achieves 80\% power at $n = 14$, confirming it is well-powered at $n = 12$. However, the VAE positive trend ($\rho = +0.33$) would require $n > 60$ for 80\% power---far beyond our matched set. The VAE result should be interpreted as a positive trend, with the method difference ($\Delta\rho$) carrying the statistical weight. Spatial coordinate matching in Allen CCF space could expand the matched set but introduces many-to-one matching ambiguities.

\textbf{IIA vacuity.} The structured VAE's IIA is potentially vacuous (\S\ref{sec:vacuity}). The shuffled-label control ($\S$5.4) resolves this: the VAE learns real structure (real IIA $> $ shuffled IIA in all 24 regions, $p = 6 \times 10^{-8}$), and the linear ablation shows the structured split, not nonlinearity, drives the improvement.

\textbf{Temporal resolution.} Sliding-window IIA analysis ($\S$5.8) shows pre-response emergence of the choice subspace in 100\% of regions with detectable onset, confirming a causal decision signal rather than post-hoc motor artifact.

\textbf{Dimensionality confound with neuron count.} The power-law exponent $\alpha$ is strongly correlated with neuron count ($\rho = -0.91$, $p < 10^{-100}$; Appendix~\ref{app:alpha_bias}), consistent with finite-sample bias identified by \citet{chun2025dimensionality}. All $\alpha$-based claims in this paper should be interpreted with this confound in mind; partial correlations controlling for neuron count are reported alongside raw correlations.

\textbf{CKA bias.} Standard (biased) CKA inflates similarity estimates in the low-trial, high-neuron regime characteristic of our data \citep{murphy2024biased}. Debiased CKA (Appendix~\ref{app:debiased_cka}) reduces the CKA--Procrustes correlation from $\rho = 0.85$ to $\rho = 0.45$, confirming the inflation but also confirming that the metric dissociation persists after correction.

\textbf{Task specificity.} All results are for binary visual choice. Other task types (navigation, working memory) may show different patterns.

\section{Conclusion}

We establish two geometric dissociations in neural population codes. First, linear and nonlinear similarity metrics anti-correlate ($\rho = -0.85$), mediated by effective dimensionality; this effect replicates on an independent multi-laboratory dataset ($\rho = -0.94$, IBL Brain-Wide Map, 139 mice). Second, linear subspace methods systematically underestimate causal structure by 3.4$\times$ and are \emph{anti-correlated} with optogenetic ground-truth causal importance ($\rho = -0.73$). These dissociations converge: the brain regions where linear methods fail most severely are precisely the regions where causal signals are strongest. Structured nonlinear methods recover these signals, producing subspaces that generalize across four task variables, disentangle engagement from choice, and better predict optogenetic silencing effects. However, the IIA metric itself is potentially vacuous for nonlinear methods, motivating external validation as the gold standard for subspace quality assessment.

\bibliographystyle{tmlr}
\bibliography{references}

\appendix
\section{Sequential Discovery and Multiple Comparisons}
\label{app:sequential}

\begin{table}[h]
\centering
\small
\begin{tabular}{@{}clcccc@{}}
\toprule
\textbf{Stage} & \textbf{Test} & \textbf{$\rho$} & \textbf{$p$} & \textbf{$N_k$} & \textbf{Survives} \\
\midrule
\multicolumn{6}{l}{\emph{Stage 1: Core metric comparison}} \\
1 & CKA--Procrustes anti-correlation & $-0.85$ & $\approx 0$ & 1 & \checkmark \\
\midrule
\multicolumn{6}{l}{\emph{Stage 2: Mechanism identification}} \\
2 & Dimensionality mediates (partial $r$) & $+0.44$ & $< 10^{-6}$ & 2 & \checkmark \\
\midrule
\multicolumn{6}{l}{\emph{Stage 3: Causal subspace discovery}} \\
3 & VAE $>$ LDA IIA (73/73 regions) & --- & $5.7 \times 10^{-14}$ & 3 & \checkmark \\
4 & Silencing vs.\ LDA IIA & $-0.73$ & $0.01$ & 4 & \checkmark \\
5 & $\Delta\rho$ (VAE $-$ LDA) vs.\ silencing & $+1.06$ & $< 0.005$ & 5 & \checkmark \\
\midrule
\multicolumn{6}{l}{\emph{Stage 4: Generalization}} \\
6 & Multi-task VAE wins (4 variables) & --- & $< 10^{-13}$ & 6 & \checkmark \\
7 & Engagement-choice orthogonality & $d_G = 1.82$ & --- & 7 & \checkmark \\
\midrule
\multicolumn{6}{l}{\emph{Stage 5: IIA vacuity and controls}} \\
8 & VAE random IIA $\approx$ MLP random IIA & $0.70$ & $p = 0.38$ (no diff) & 8 & (diagnostic) \\
\bottomrule
\end{tabular}
\caption{Sequential discovery table. All seven core claims (stages 1--4) survive Holm correction. Stage 5 is a diagnostic test, not a directional hypothesis.}
\label{tab:sequential}
\end{table}

\section{Robustness Controls}
\label{app:robustness}

\begin{table}[h]
\centering
\small
\caption{Robustness of the CKA--Procrustes anti-correlation.}
\label{tab:robustness}
\begin{tabular}{@{}lccl@{}}
\toprule
\textbf{Control} & \textbf{$\rho$} & \textbf{$p$} & \textbf{$n$} \\
\midrule
Raw (full data) & $-0.850$ & $\approx 0$ & 1,316 pairs \\
Bootstrap 95\% CI & $[-0.867, -0.830]$ & --- & 10,000 resamples \\
Leave-one-mouse-out (range) & $[-0.872, -0.835]$ & all sig. & 10 mice \\
Leave-one-region-out (range) & $[-0.863, -0.841]$ & all sig. & 50 regions \\
Neuron sub-sampling & $-0.814$ & $< 10^{-7}$ & 29 pairs \\
Spectral denoising (MP) & $-0.889$ & $\approx 0$ & 1,316 pairs \\
Trial-label shuffling & $-0.395$ & $< 10^{-50}$ & 1,316 pairs \\
\midrule
\multicolumn{4}{l}{\emph{Partial correlations (Pearson $r$)}} \\
Controlling neuron count & $+0.134$ & --- & sign reversal \\
Controlling $\alpha$ & $+0.437$ & --- & sign reversal \\
Controlling both & $+0.196$ & --- & sign reversal \\
\bottomrule
\end{tabular}
\end{table}

\section{Metric Family Comparison}
\label{app:metrics}

\begin{table}[h]
\centering
\small
\caption{Correlation of each metric family with dimensionality ($\alpha$) and with CKA/Procrustes.}
\label{tab:metrics}
\begin{tabular}{@{}lccccc@{}}
\toprule
\textbf{Metric} & \textbf{vs.\ $\alpha$} & \textbf{$p$} & \textbf{vs.\ CKA} & \textbf{vs.\ Procrustes} \\
\midrule
CKA (linear) & $-0.68$ & $< 10^{-4}$ & --- & $-0.85$ \\
UMAP Procrustes & $+0.68$ & $< 10^{-4}$ & $-0.85$ & --- \\
Sliced Wasserstein & $+0.02$ & $0.47$ & $+0.10$ & $-0.09$ \\
SPD Riemannian & $+0.44$ & $< 10^{-4}$ & --- & --- \\
Fisher-Rao & $+0.40$ & $< 10^{-4}$ & --- & --- \\
\bottomrule
\end{tabular}
\end{table}

\section{Evidence-Choice Subspace Alignment}
\label{app:alignment}

\begin{table}[h]
\centering
\small
\caption{Top and bottom 5 regions by evidence-choice subspace alignment.}
\label{tab:routing}
\begin{tabular}{@{}lccc@{}}
\toprule
\textbf{Region} & \textbf{$d_G(V_{\text{ev}}, V_{\text{ch}})$} & \textbf{$\alpha$} & \textbf{Interpretation} \\
\midrule
\multicolumn{4}{l}{\emph{Highest alignment (shared subspace)}} \\
OT & 0.197 & 104.2 & evidence-choice routing \\
IC & 0.409 & 13.1 & sensorimotor integration \\
MEA & 0.453 & 3.0 & evidence-choice routing \\
LH & 0.471 & 3.9 & evidence-choice routing \\
SCm & 0.486 & 15.6 & motor output \\
\midrule
\multicolumn{4}{l}{\emph{Lowest alignment (separated subspaces)}} \\
VAL & 1.163 & 1.7 & thalamic relay \\
LSc & 1.216 & 2.4 & basal ganglia \\
COA & 1.313 & 0.7 & amygdala \\
CL & 1.381 & 27.7 & thalamic relay \\
EPd & 1.397 & 43.6 & amygdala/piriform \\
\bottomrule
\end{tabular}
\end{table}

\section{Topological and Dynamical Invariants}
\label{app:topo_dyn}

\textbf{Persistent homology.} Only 9/73 regions have $\beta_1 > 0$. Grassmannian vs.\ topological Wasserstein $H_1$ distance: $\rho = 0.25$, $p = 0.52$ ($n = 9$).

\textbf{Rotational dynamics.} 72/73 regions show rotation above null. Mean strength $= 0.053 \pm 0.060$; range $0.008$--$0.40$.

\textbf{Spectral universality.} Singular value spectra: mean $r = 0.94 \pm 0.06$ (1,316 pairs). Temporal modes: mean cosine $= 0.09 \pm 0.07$.

\section{Grassmannian Distance and Functional Hierarchy}
\label{app:gdm_hierarchy}

The GDM does not recover traditional functional groupings. We assigned 39 of 73 regions to broad functional groups (visual, motor, somatosensory, auditory, thalamic, hippocampal, prefrontal) based on the Allen Brain Atlas, yielding 97 within-group and 644 between-group pairs. Within-group Grassmannian distances are not smaller than between-group distances for either LDA (Mann-Whitney $U = 34{,}509$, $p = 0.97$) or the VAE ($U = 35{,}465$, $p = 0.99$). Spectral clustering of the GDM into $k = 7$ groups yields weak agreement with functional labels (LDA ARI $= 0.07$, NMI $= 0.18$; VAE ARI $= 0.07$, NMI $= 0.20$). This null result suggests that choice subspace geometry is organized along axes orthogonal to the classical sensory--association--motor hierarchy---consistent with the view that the choice signal is routed through a distributed network whose geometry does not respect anatomical proximity or functional category.

\section{Tucker Decomposition}
\label{app:tucker}

Tucker decomposition on (trials $\times$ neurons $\times$ time) tensors: time factors are most conserved (mean similarity $0.174$), followed by trial ($0.106$) and neuron ($0.067$). Tucker $R^2$ correlates modestly with $\alpha$ ($\rho = 0.275$).

\section{Debiased CKA}
\label{app:debiased_cka}

\citet{murphy2024biased} showed that the standard CKA estimator (using centered kernel alignment with the biased HSIC) produces inflated similarity scores when the sample-to-feature ratio is small---exactly our regime (250 trials, 15--40 neurons per region, unequal region sizes across pairs). We re-ran the main CKA--Procrustes analysis using the unbiased HSIC$_1$ estimator of \citet{song2012feature}:

\begin{table}[h]
\centering
\small
\caption{Biased vs.\ debiased CKA correlation with UMAP Procrustes similarity across 1,316 cross-session pairs.}
\label{tab:debiased_cka}
\begin{tabular}{@{}lccc@{}}
\toprule
\textbf{Estimator} & \textbf{Spearman $\rho$} & \textbf{$p$} & \textbf{Partial $\rho$ (ctrl.\ $n_{\text{neurons}}$)} \\
\midrule
Biased CKA & $0.850$ & $\approx 0$ & $0.856$ \\
Debiased CKA & $0.445$ & $4.2 \times 10^{-65}$ & $0.342$ \\
\bottomrule
\end{tabular}
\end{table}

Debiasing reduces the CKA--Procrustes correlation by $\Delta\rho = -0.40$, confirming that the biased estimator inflates apparent similarity. Critically, the debiased correlation remains highly significant ($p < 10^{-60}$), meaning the metric dissociation is real but weaker than the biased estimator suggests. The partial correlation controlling for neuron count drops further to $0.34$, indicating that neuron-count differences between regions contribute additional apparent correlation beyond the CKA inflation.

\section{Power-Law Alpha and Neuron Count Confound}
\label{app:alpha_bias}

\citet{chun2025dimensionality} showed that finite-sample bias in PCA eigenvalue estimation inflates power-law exponent estimates for small neuron counts. We tested whether our $\alpha$ estimates are confounded by recording yield:

\begin{table}[h]
\centering
\small
\caption{Confound analysis: power-law exponent vs.\ recording parameters (325 region-sessions, 283 with valid $\alpha$).}
\label{tab:alpha_confound}
\begin{tabular}{@{}lcc@{}}
\toprule
\textbf{Covariate} & \textbf{Spearman $\rho$ with $\alpha$} & \textbf{$p$} \\
\midrule
$n_{\text{neurons}}$ & $-0.908$ & $8.6 \times 10^{-103}$ \\
$n_{\text{trials}}$ & $+0.106$ & $0.083$ \\
\bottomrule
\end{tabular}
\end{table}

The neuron count confound is massive ($|\rho| = 0.91$): regions with more recorded neurons systematically show lower $\alpha$ (higher effective dimensionality). This likely reflects both (1) finite-sample bias in eigenvalue estimation and (2) a genuine biological correlation---larger regions may have higher intrinsic dimensionality. Trial count is not confounded ($p = 0.08$). All $\alpha$-based claims in this paper should be interpreted with this caveat; where possible, we report partial correlations controlling for neuron count.

Bootstrap confidence intervals for $\alpha$ estimates (200 resamples per region-session) show wide variability: mean bootstrap $\text{SD} > 1$ for most regions, consistent with the high sensitivity to sampling identified by \citet{chun2025dimensionality}.

\section{Shuffled-Label and Linear Ablation Controls}
\label{app:shuffled_ablation}

\textbf{Shuffled-label control (Exp62).} We trained the structured VAE with randomly shuffled choice labels on real neural data (5 shuffles per region, 24 regions with sufficient data). Results:

\begin{table}[h]
\centering
\small
\caption{IIA under real vs.\ shuffled labels (24 regions).}
\label{tab:shuffled}
\begin{tabular}{@{}lccc@{}}
\toprule
\textbf{Condition} & \textbf{Mean IIA} & \textbf{$n$} & \textbf{Interpretation} \\
\midrule
Real labels & $0.625 \pm 0.089$ & 24 & Genuine structure \\
Shuffled choice & $0.430 \pm 0.060$ & 24 & Encoder flexibility \\
Shuffled evidence & $0.424 \pm 0.043$ & 24 & No evidence structure \\
\bottomrule
\end{tabular}
\end{table}

Real labels produce significantly higher IIA in all 24/24 regions (Wilcoxon $W = 300$, $p = 6.0 \times 10^{-8}$). The gap between real and shuffled ($\Delta = 0.195$) is the genuine signal; the shuffled baseline ($0.43$) is the encoder flexibility floor (cf.\ the $0.70$ noise-fitting on random Gaussian data in $\S$\ref{sec:vacuity}---the difference reflects that shuffled labels on real data are harder to fit than labels on i.i.d.\ Gaussian noise).

\textbf{Linear ablation (Exp63).} Nonlinearity contributes nothing beyond the structured split:

\begin{table}[h]
\centering
\small
\caption{Ablation: nonlinear vs.\ linear encoder with structured split (24 regions).}
\label{tab:linear_ablation}
\begin{tabular}{@{}lcc@{}}
\toprule
\textbf{Effect} & \textbf{Mean $\Delta$IIA} & \textbf{Wilcoxon $p$} \\
\midrule
Nonlinearity (nonlinear $-$ linear structured) & $+0.001$ & $0.65$ \\
\bottomrule
\end{tabular}
\end{table}

\section{Optogenetic Power Analysis}
\label{app:power}

A power analysis for the optogenetic Spearman correlation clarifies the interpretation:

\begin{table}[h]
\centering
\small
\caption{Statistical power for optogenetic correlations at $\alpha = 0.05$ (10,000 simulations per sample size).}
\label{tab:power}
\begin{tabular}{@{}lccl@{}}
\toprule
\textbf{Method} & \textbf{$\rho$} & \textbf{$n$ for 80\% power} & \textbf{Power at $n = 12$} \\
\midrule
LDA & $-0.73$ & 14 & 0.74 \\
VAE & $+0.33$ & $> 60$ & 0.16 \\
\bottomrule
\end{tabular}
\end{table}

The LDA anti-correlation is nearly adequately powered at $n = 12$ (74\%) and crosses 80\% at $n = 14$. The VAE positive trend requires $n > 60$ for adequate power---well beyond our matched set. This motivates emphasis on the method difference ($\Delta\rho$) rather than the VAE result alone, and motivates future work expanding the matched region set via Allen CCF coordinate matching.

\section{Per-Mouse Optogenetic Validation}
\label{app:per_mouse}

Results from per-mouse silencing analysis are reported in the supplementary data files (Exp66).

\section{pi-VAE Hybrid Ablation}
\label{app:pivae}

We compared four encoder architectures across 73 regions to test whether adding a label-conditioned prior (pi-VAE; \citealp{zhou2020pivae}) improves upon the structured split:

\begin{table}[h]
\centering
\small
\caption{Model comparison: mean IIA across 73 regions.}
\label{tab:pivae}
\begin{tabular}{@{}lc@{}}
\toprule
\textbf{Model} & \textbf{Mean IIA} \\
\midrule
Structured VAE & $\mathbf{0.941}$ \\
pi-structured-VAE hybrid & $0.036$ \\
pi-VAE (no split) & $0.026$ \\
LDA & $0.000$ \\
\bottomrule
\end{tabular}
\end{table}

The structured VAE dramatically outperforms all alternatives. The pi-VAE and hybrid both achieve near-zero IIA, performing far worse than the structured VAE alone (Wilcoxon $p = 9.4 \times 10^{-14}$, 0/73 wins for the hybrid). The pi-VAE's label-conditioned prior does not improve and in fact disrupts the structured split when combined. This suggests that the structured split's inductive bias (separate $\mathbf{z}_{\text{choice}}$ and $\mathbf{z}_{\text{other}}$ dimensions) is the critical architectural choice, and adding identifiability constraints via the prior is unnecessary or even harmful in this regime.

\section{Cross-Region Activation Patching}
\label{app:patching}

We performed biological path patching across all simultaneously recorded region pairs (1,438 directed pairs from 39 sessions). For each pair $A \to B$, we replaced region $B$'s choice-subspace projection with region $A$'s and measured whether a downstream classifier flipped its choice prediction (patching IIA).

\begin{table}[h]
\centering
\small
\caption{Top hub regions by mean outgoing patching IIA (73 regions).}
\label{tab:hubs}
\begin{tabular}{@{}lcccl@{}}
\toprule
\textbf{Region} & \textbf{Out IIA} & \textbf{In IIA} & \textbf{$n$ pairs} & \textbf{Asymmetry} \\
\midrule
MG (medial geniculate) & 0.651 & 0.594 & 13 & $+0.057$ \\
MD (mediodorsal thal.) & 0.629 & 0.593 & 18 & $+0.036$ \\
ACB (nucleus accumbens) & 0.573 & 0.666 & 22 & $-0.094$ \\
SUB (subiculum) & 0.536 & 0.495 & 32 & $+0.041$ \\
NB (basal forebrain) & 0.500 & 0.489 & 8 & $+0.011$ \\
VISa (anterior visual) & 0.491 & 0.580 & 27 & $-0.089$ \\
root (unassigned) & 0.430 & 0.388 & 59 & $+0.042$ \\
ZI (zona incerta) & 0.408 & 0.405 & 25 & $+0.003$ \\
\bottomrule
\end{tabular}
\end{table}

The asymmetry between outgoing and incoming patching IIA reveals directional information flow. Thalamic regions (MG, MD) show positive asymmetry---they are causal sources whose choice signal propagates to cortex. Conversely, ACB and VISa show negative asymmetry (net receivers). The top patching pairs (CA3$\to$LS, SUB$\to$LS) achieve IIA $= 1.0$, indicating perfect choice-signal transfer between hippocampal and septal regions.

\section{VAE Causal Circuits}
\label{app:circuits}

Following \citet{martinez2025vae_circuits}, we computed causal effect strength (CES), intervention specificity, and circuit modularity for the structured VAE encoder across 73 regions.

\begin{table}[h]
\centering
\small
\caption{VAE encoder causal circuit analysis (73 regions).}
\label{tab:circuits}
\begin{tabular}{@{}lc@{}}
\toprule
\textbf{Metric} & \textbf{Value} \\
\midrule
CES (choice dimensions) & $0.077$ \\
CES (other dimensions) & $0.028$ \\
Choice/other CES ratio & $\mathbf{2.73}$ \\
Random-split CES ratio & $1.005 \pm 0.024$ \\
Lift over random split & $2.72\times$ \\
Structured modularity & $0.859$ \\
Unstructured modularity & $0.926$ \\
\bottomrule
\end{tabular}
\end{table}

The choice dimensions carry $2.73\times$ more causal effect strength than the other dimensions---and $2.72\times$ more than a random split of dimensions, confirming that the structured VAE develops specialized choice-encoding channels rather than distributing information uniformly. The unstructured VAE achieves slightly higher modularity ($0.926$ vs.\ $0.859$), likely because it has no forced split to constrain its internal organization.

\section{UMAP Stochasticity Robustness}
\label{app:umap}

Because UMAP is stochastic (different random seeds yield different embeddings), we verified that the CKA--Procrustes correlation is not an artifact of a particular UMAP realization. We swept 20 random seeds $\times$ 9 hyperparameter configurations (3 neighbor counts $\times$ 3 minimum distances) across 24 regions.

\begin{table}[h]
\centering
\small
\caption{Spearman $\rho$ (CKA vs.\ Procrustes) across 20 UMAP seeds per configuration.}
\label{tab:umap}
\begin{tabular}{@{}ccccc@{}}
\toprule
\textbf{$n_{\text{neighbors}}$} & \textbf{min\_dist} & \textbf{Mean $\rho$} & \textbf{SD} & \textbf{Range} \\
\midrule
5 & 0.0 & 0.888 & 0.014 & [0.848, 0.910] \\
5 & 0.1 & 0.898 & 0.010 & [0.876, 0.914] \\
5 & 0.5 & 0.900 & 0.011 & [0.882, 0.922] \\
15 & 0.0 & 0.858 & 0.016 & [0.828, 0.885] \\
15 & 0.1 & 0.850 & 0.012 & [0.833, 0.877] \\
15 & 0.5 & 0.820 & 0.021 & [0.775, 0.849] \\
30 & 0.0 & 0.875 & 0.008 & [0.859, 0.891] \\
30 & 0.1 & 0.867 & 0.010 & [0.847, 0.886] \\
30 & 0.5 & 0.826 & 0.020 & [0.784, 0.865] \\
\bottomrule
\end{tabular}
\end{table}

Across all 180 runs (20 seeds $\times$ 9 configs), the CKA--Procrustes Spearman $\rho$ ranges from $0.775$ to $0.922$ (grand mean $0.864$, SD across configs $= 0.028$). No configuration produces a $\rho$ below $0.77$. The anti-correlation is robust to UMAP stochasticity and hyperparameter choice, with larger $n_{\text{neighbors}}$ and larger min\_dist slightly reducing the correlation (as expected, since these produce more global embeddings that compress local structure).

\section{pi-VAE Identifiability Failure Analysis}
\label{app:ivae-verification}

The pi-VAE hybrid (Appendix~\ref{app:pivae}) achieves IIA $= 0.036$, far below the structured VAE's $0.941$. We measure the iVAE identifiability conditions \citep{khemakhem2020variational} to explain this mechanistically. The key requirement is that the label-conditional prior means span a space of dimension $\geq d_{\text{causal}}$. With binary choice ($n_{\text{classes}} = 2$) and $d_{\text{causal}} = 3$, the rank condition provably fails: two class means produce at most rank~1.

\begin{table}[h]
\centering
\caption{iVAE identifiability metrics across 73 brain regions ($n_{\text{classes}} = 2$, $d_{\text{causal}} = 3$). The rank condition is never satisfied for any model variant.}
\label{tab:ivae-verification}
\small
\begin{tabular}{lccccc}
\toprule
\textbf{Model} & \textbf{Max $\rho$} & \textbf{Eff.\ rank} & \textbf{Recon.\ MSE} & \textbf{$z$ sep.} \\
\midrule
Structured VAE       & 0.767 & 2.00 & 0.0022  & 4.75 \\
pi-VAE               & 0.063 & 1.98 & 1.55    & 0.99 \\
pi-Structured VAE    & 0.066 & 1.90 & 0.0037  & 0.51 \\
\bottomrule
\end{tabular}
\end{table}

The structured VAE achieves Spearman $\rho = 0.77$ between $z_{\text{choice}}$ and true labels, with well-separated class means ($z$-separation $= 4.75$) and low reconstruction error. The pi-VAE variants show near-zero correlation ($\rho \approx 0.06$), and pi-VAE's reconstruction MSE is $700\times$ worse ($1.55$ vs.\ $0.0022$). The label-conditional prior forces the latent space toward two point masses, destroying the decoder's ability to reconstruct neural activity. The pi-structured hybrid rescues reconstruction (MSE $= 0.0037$) via its classification head, but the prior still collapses the latent structure ($z$-separation $= 0.51$).

The effective rank is $\leq 2$ for all models (matching the number of classes), never reaching the $d_{\text{causal}} = 3$ required by Theorem~1 of \citet{khemakhem2020variational}. This confirms that the pi-VAE failure is not an optimization accident but a fundamental mismatch: binary choice data cannot satisfy the identifiability conditions when the causal subspace is higher-dimensional than the number of label classes.

\section{CD-NOD Causal Discovery Over Brain Regions}
\label{app:cdnod}

We apply CD-NOD \citep{huang2020causal} to discover directed causal structure between brain regions from the VAE choice subspace projections. CD-NOD exploits distribution shifts across recording sessions (different mice, different days) to orient edges that purely observational methods cannot. For each region and session, we train a structured VAE and extract the first $z_{\text{choice}}$ component. The data matrix has shape (trials $\times$ regions), with session index as the context variable $c$.

Sparse session overlap in the Steinmetz dataset limits the analysis to 3 co-occurring regions (CA1, DG, root) across 7 sessions (1,648 trials). We compare three methods:

\begin{table}[h]
\centering
\caption{Causal discovery results on VAE choice subspace projections (3 regions, 7 sessions, 1,648 trials).}
\label{tab:cdnod}
\small
\begin{tabular}{lcl}
\toprule
\textbf{Method} & \textbf{Directed edges} & \textbf{Key finding} \\
\midrule
CD-NOD         & 3 & DG $\to$ CA1, DG $\to$ root, root $\to$ CA1 \\
PC algorithm   & 0 (3 undirected) & Finds connections but cannot orient \\
DirectLiNGAM   & 3 & CA1 $\to$ DG, root $\to$ DG (reversed) \\
\bottomrule
\end{tabular}
\end{table}

CD-NOD recovers the anatomically correct DG $\to$ CA1 direction (dentate gyrus projects to CA1 via CA3 in the hippocampal trisynaptic circuit), while PC leaves all edges undirected and DirectLiNGAM reverses the DG--CA1 edge. CD-NOD also identifies ``root'' as a changing module --- its causal mechanism varies across sessions, consistent with the heterogeneous composition of this aggregate region. While the 3-region scope limits the strength of this analysis, it demonstrates that session heterogeneity provides genuine causal orientation information on VAE subspace projections. Denser recordings (e.g., IBL repeated-site data) would enable the same pipeline at brain-wide scale.

\end{document}
